\chapter{ปรากฏการณ์ความร้อนและการถ่ายโอนความร้อน}
\label{ch:ch09_heat_transfer}
\index{การถ่ายโอนความร้อน}

\begin{conceptbox}[title=วัตถุประสงค์การเรียนรู้ประจำบท]
เมื่อสิ้นสุดการศึกษาบทเรียนนี้ นักศึกษาสามารถ
\begin{enumerate}
    \item แยกความแตกต่างระหว่างความร้อนและอุณหภูมิ พร้อมอธิบายความหมายเชิงจุลภาคของทั้งสองปริมาณได้อย่างถูกต้อง
    \item แปลงหน่วยอุณหภูมิระหว่างมาตราเซลเซียส เคลวิน และฟาเรนไฮต์ และคำนวณการขยายตัวเนื่องจากความร้อนของวัสดุ
    \item คำนวณปริมาณความร้อนที่เกี่ยวข้องกับการเปลี่ยนอุณหภูมิของสารโดยใช้สมการ Q = mc$\Delta$T
    \item คำนวณปริมาณความร้อนแฝงที่เกี่ยวข้องกับการเปลี่ยนสถานะของสารโดยใช้สมการ Q = mL
    \item อธิบายกลไกการถ่ายเทความร้อนทั้งสามแบบ ได้แก่ การนำ การพา และการแผ่รังสี พร้อมยกตัวอย่างประกอบ
\end{enumerate}
\end{conceptbox}

\section{ทฤษฎีและหลักการพื้นฐาน}

\begin{bioappbox}
\textbf{การควบคุมอุณหภูมิในสิ่งมีชีวิตและการฆ่าเชื้อด้วยความร้อน}

การถ่ายโอนความร้อนทั้ง 3 รูปแบบ ได้แก่ การนำความร้อน (Conduction) การพาความร้อน (Convection) และการแผ่รังสีความร้อน (Radiation) มีบทบาทสำคัญต่อสรีรวิทยาของสิ่งมีชีวิตในการรักษาอุณหภูมิร่างกายให้อยู่ในภาวะสมดุล (Thermoregulation) นอกจากนี้ ในทางจุลชีววิทยา หลักการความร้อนแฝงของการกลายเป็นไอ (Latent Heat of Vaporization) ถูกนำมาใช้ในหม้อนึ่งความดันไอ (Autoclave) ที่อุณหภูมิ $121^\circ\text{C}$ ความดัน 15 psi เพื่อทำลายสปอร์ของแบคทีเรียได้อย่างสมบูรณ์
\end{bioappbox}

\begin{labbox}
1. ตรวจสอบตำแหน่งศูนย์ (Zero Error) ของเครื่องมือวัดทุกชนิดก่อนเริ่มบันทึกข้อมูล
2. ทำการวัดซ้ำอย่างน้อย 3 ถึง 5 ครั้งในแต่ละจุดทดลองเพื่อคำนวณหาค่าเฉลี่ยและค่าเบี่ยงเบนมาตรฐาน
3. บันทึกผลการวัดพร้อมระบุหน่วยเอสไอ (SI Units) และค่าความคลาดเคลื่อนของการวัดเสมอ
\end{labbox}

\section{ตัวอย่างโจทย์และการวิเคราะห์เชิงฟิสิกส์}

\begin{examplebox}
การแผ่รังสีความร้อน (radiation) เป็นการถ่ายเทความร้อนในรูปของคลื่นแม่เหล็กไฟฟ้า โดยเฉพาะในช่วงอินฟราเรด ซึ่งไม่จำเป็นต้องอาศัยตัวกลางใด ๆ ในการถ่ายเท สามารถเดินทางผ่านสุญญากาศได้ ดังตัวอย่างที่ชัดเจนที่สุดคือพลังงานความร้อนจากดวงอาทิตย์ที่เดินทางผ่านอวกาศมายังโลก วัตถุทุกชนิดที่มีอุณหภูมิสูงกว่าศูนย์สัมบูรณ์จะแผ่รังสีความร้อนออกมาตลอดเวลา โดยอัตราการแผ่รังสีความร้อนของวัตถุแปรผันตรงกับอุณหภูมิสัมบูรณ์ยกกำลังสี่ ตามกฎของสเตฟาน-โบลต์ซมันน์ (Stefan–Boltzmann law) วัตถุที่มีผิวสีเข้มและหยาบจะดูดกลืนและแผ่รังสีความร้อนได้ดีกว่าวัตถุที่มีผิวมันวาวสีอ่อน ซึ่งเป็นหลักการที่นำไปใช้ในการออกแบบเสื้อผ้า อุปกรณ์กันความร้อน และแผงรับพลังงานแสงอาทิตย์

ตารางสรุปสูตรสำคัญประจำบท

ปริมาณ / กฎ

สมการ
\end{examplebox}

\section{แบบฝึกหัดทบทวนและคำถามท้ายบท}

\begin{enumerate}
    \item จงอธิบายความแตกต่างระหว่างความร้อนกับอุณหภูมิ พร้อมยกตัวอย่างประกอบ
    \item จงแปลงอุณหภูมิ 37 องศาเซลเซียส (อุณหภูมิร่างกายมนุษย์โดยประมาณ) เป็นหน่วยเคลวินและฟาเรนไฮต์
    \item จงแปลงอุณหภูมิ 98.6 องศาฟาเรนไฮต์ เป็นหน่วยองศาเซลเซียส
    \item เส้นลวดเหล็กยาว 10.0 เมตร ที่อุณหภูมิ 20 องศาเซลเซียส จะมีความยาวเปลี่ยนแปลงไปเท่าใดเมื่ออุณหภูมิเพิ่มขึ้นเป็น 45 องศาเซลเซียส (กำหนด $\alpha$ เหล็ก = 12$\times$10$^-$⁶ $^\circ$C$^-$$^1$)
    \item ทำไมสะพานเหล็กและทางรถไฟจึงต้องมีรอยต่อขยายตัว (expansion joint) จงอธิบายด้วยหลักการทางฟิสิกส์
    \item จงคำนวณปริมาณความร้อนที่ต้องใช้ในการทำให้แท่งอะลูมิเนียมมวล 0.80 กิโลกรัม มีอุณหภูมิเพิ่มขึ้นจาก 20 องศาเซลเซียส เป็น 100 องศาเซลเซียส (กำหนด c อะลูมิเนียม = 900 J/(kg·K))
\end{enumerate}

% สิ้นสุดบทเรียน
